\newpage
\section{问题分析}

\subsection{问题一分析}
交通流数据具有显著的\textbf{时空相关性}——相邻路口的流量相互影响，同一路口不同时段也存在自相关性。我们需要从数据中提取时空特征，为后续预测和控制提供输入。

\subsection{问题二分析}
多路口协同控制是\textbf{大规模组合优化}问题。相位差优化变量是连续量（每个路口的延迟时间），约束包括绿波带宽、最小红黄灯时间等。我们采用改进的遗传算法求解。

\subsection{问题三分析}
强化学习控制的关键是\textbf{状态-动作-奖励}设计：状态包含各方向排队长度和相位信息，动作为相位切换和绿灯时长调整，奖励为负总延误时间。Q-learning适合离散动作空间，DQN适合连续动作空间。
